\scanpage{29}
\mspage{2}

jouant le rôle de ``pierres de construction'' élémentaires. Ceci amène à introduire
\[
(5)\qquad \unsure{\mathcal V_{Y,X}^*|V_{Y,X}^*}\overset{\mathrm{déf}}=\unsure{\mathcal V_{Y,X}|V_{Y,X}}\cap X^*=\unsure{\mathcal V_{Y,X}|V_{Y,X}}\setminus Y
\]
\[
(X^*=X-Y),
\]
et le diagramme \deleted{\unsure{un}}

(6)
\diagram{\begin{tikzcd}[ampersand replacement=\&,column sep=large]
\unsure{\mathcal V_{Y,X}^*|V_{Y,X}^*}\arrow[r,hook,"\substack{\text{bouchage de trous}\\\text{le long de }Y}"]\arrow[d]\&\unsure{\mathcal V_{Y,X}|V_{Y,X}}\arrow[d]\\
X^*\arrow[r]\&X
\end{tikzcd}}

Je vais regarder (indépendamment de toute hypothèse de locale trivialité normale, ou de lissité de $X^*$ et $Y$) $X$ comme obtenu par \uline{recollement} de $\unsure{\mathcal V_{Y,X}|V_{Y,X}}$ et de $X^*$ suivant leur intersection $\unsure{\mathcal V_{Y,X}^*|V_{Y,X}^*}$, ou (en termes plus techniques) comme la somme amalgamée, ou comme \deleted{\unsure{espace}} espace-topos $\varinjlim$ du diagramme

(7)
\diagram{\begin{tikzcd}[ampersand replacement=\&]
\unsure{\mathcal V_{Y,X}^*|V_{Y,X}^*}\arrow[r,hook]\arrow[d]\&\unsure{\mathcal V_{Y,X}|V_{Y,X}}\\
X^*\&
\end{tikzcd}}

\margin{28}
